\section{USAMO 2020}
        \begin{enumerate}
        	\item (\textit{USAMO 2020}) Let $ABC$ be a fixed acute triangle inscribed in a circle $\omega$ with center $O$. A variable point $X$ is chosen on minor arc $AB$ of $\omega$, and segments $CX$ and $AB$ meet at $D$. Denote by $O_1$ and $O_2$ the circumcenters of triangles $ADX$ and $BDX$, respectively. Determine all points $X$ for which the area of triangle $OO_1O_2$ is minimized.


 



 



\textbf{Solution}


                    \begin{figure}[H]
                        \centering
                        \includegraphics[width=0.5\textwidth]{Solution/USAMO_I_2020_1_Solution_Figure_1.png}
                    \end{figure}
                    Let $E$ be midpoint $AD.$ Let $F$ be midpoint $BD \implies$
 \[EF = ED + FD = \frac {AD}{2} + \frac {BD}{2} = \frac {AB}{2}.\]$E$ and $F$ are the bases of perpendiculars dropped from $O_1$ and $O_2,$ respectively.

 Therefore  $O_1O_2 \ge EF =  \frac {AB}{2}.$

 
 \[CX \perp  O_1O_2, AX \perp  O_1O \implies \angle O O_1O_2 = \angle AXC\]$\angle AXC = \angle ABC (AXBC$ is cyclic) $\implies \angle O O_1O_2 = \angle ABC.$

 Similarly $\angle BAC = \angle O O_2 O_1 \implies  \triangle ABC \sim \triangle O_2 O_1O.$ 

 The area of $\triangle OO_1O_2$ is minimized if $CX \perp AB$ because
 \[\frac {[OO_1O_2]} {[ABC]} = \left(\frac {O_1 O_2} {AB}\right)^2 \ge \left(\frac {EF} {AB}\right)^2 = \frac {1}{4}.\]\textbf{vladimir.shelomovskii@gmail.com, vvsss}

 


	\item (\textit{USAMO 2020}) An empty $2020 \times 2020 \times 2020$ cube is given, and a $2020 \times 2020$ grid of square unit cells is drawn on each of its six faces. A beam is a $1 \times 1 \times 2020$ rectangular prism. Several beams are placed inside the cube subject to the following conditions:


 
                \begin{itemize}
                    \item The two $1 \times 1$ faces of each beam coincide with unit cells lying on opposite faces of the cube. (Hence, there are $3 \cdot 2020^2$ possible positions for a beam.)
\item No two beams have intersecting interiors.
\item The interiors of each of the four $1 \times 2020$ faces of each beam touch either a face of the cube or the interior of the face of another beam.
                \end{itemize}


                What is the smallest positive number of beams that can be placed to satisfy these conditions?


 



\textbf{Solution}

Take one vertex of the cube as origin and establish 3D coordinates along the cube's edges.

 Define a beam as $x-dir$ if its long edge is parallel to x-axis. Similarly for $y-dir$ and $z-dir$.

 Define a beam's location as (direction, ($1 \times 1$ face's location in 2D coordinate).

 For example, (y, 2, 4) indicates the beam with vertex (1, 0, 3) and (2, 2020, 4)

 Apparently $x$ beam needs the other $x$ or $y$ beams to touch its $xy$ faces, $x$ or $z$ beams to touch its $xz$ faces. Similarly for $y$ and $z$ beam.

 If there are only 1-dir or 2-dirs beams, it is easy to approve that $2020 \times 2020$ is the minimal number beams.

 (for example, if only use $x-dir$ and $y-dir$ beams, all the $x-dir$ beams's xz faces can only be touched by $x-dir$ beams. In the other word, $2020 x-dir$ beams will be stacked to meet xz faces touch requirements in that xz layer)

 Consider cases with all $3-dirs$ beams.WLOG there is a $x-dir$ beam and it needs $x-dir$ or $y-dir$ beams to touch its $xy$ faces (unless it touches the cube surface). 

 And this $y-dir$ beam also needs a $x-dir$ or $y-dir$ to touch it's $xy$ faces. And so on until one which touches cube surface. So from $xy$ face perspective, it needs $2020$ beams.

 Similarly from $xz$ and $yz$ face perspective, it also needs $2020$ and $2020$ beams.

 Consider one beam has four $1 \times 2020$ faces and it can be counted twice. So there should be at least $2020 \times 3 \div 2=3030$ beams.

 Here is one solution with 3030 beams.

 $(x, 1, 1),\ (y, 1, 2),\ (z, 2, 2),$

 $\cdots ,$

 $(x, (2n+1), (2n+1)),\ (y, (2n+1), (2n+2)),\ (z, (2n+2), (2n+2)),$

 $\cdots ,$

 $(x, (2019, 2019)),\ (y, 2019, 2020),\ (z, 2020, 2020)$

 


	\item (\textit{USAMO 2020}) Let $p$ be an odd prime. An integer $x$ is called a \textit{quadratic non-residue} if $p$ does not divide $x - t^2$ for any integer $t$.


 Denote by $A$ the set of all integers $a$ such that $1 \le a < p$, and both $a$ and $4 - a$ are quadratic non-residues. Calculate the remainder when the product of the elements of $A$ is divided by $p$.


 



 



\textbf{Solution}

We will be using finite field theory and the Frobenius endomorphism freely. Also we will just give a sketch and leave out details like it's out of fashion.

 \textbf{Case 1}: $p\equiv 3\pmod{4}$. Work in $\mathbb{F}_{p^2} = \mathbb{F}_p[i]$. Suppose $a\in A$. Then, $a=-x^2$ for some $x\in\mathbb{F}_p$ and $4-a=y^2$ for some $y\in\mathbb{F}_p$, so we have 
 \[x^2+y^2 = -4.\] This factors as 
 \[(x+iy)(x-iy) = (2i)^2,\] so we have $x+iy=2it$ and $x-iy=2i/t$ for some $t\in\mathbb{F}_p[i]$, so 
 \[x=i(t+1/t)\quad\text{and}\quad y=t-1/t.\] For both of these to be in $\mathbb{F}_p$, we need $x^p=x$ and $y^p=y$, which when we work out with frobenius, tells us that $t^{p+1}+1=0$.

 Thus, if $a\in A$, we must have $a=(t+1/t)^2$ for some $t\in\mathbb{F}_p[i]$ such that $t^{p+1}+1=0$. Actually it is easy to see that if $t$ satisfies $t^{p+1}+1=0$, then both $a$ and $4-a$ are QNRs (just reverse the frobenius calculation above). Thus, 
 \[A = \{(t+1/t)^2 : t\in\mathbb{F}_p[i], t^{p+1}+1=0\}.\] Note $(t+1/t)^2 = (t'+1/t')^2$ if and only if $t'\in\{t,-t,1/t,-1/t\}$, and also $t^{p+1}+1=0\iff (t')^{p+1}+1=0$. Note that 
 \[t^{p+1}+1\mid t^{p^2}-t\] as $2(p+1)\mid p^2-1$, so $t^{p+1}+1$ has distinct roots in $\mathbb{F}_p[i]$. We see that 
 \[t^{p+1}+1=(t^{\frac{p+1}{2}}+i)(t^{\frac{p+1}{2}}-i),\] and $s$ a root of $t^{\frac{p+1}{2}}-i$ if and only if $1/s$ is a root. Thus, it is easy to see that 
 \[A = \{(t+1/t)^2 : t\in\mathbb{F}_p[i], t^{\frac{p+1}{2}}-i=0\}.\] In this parameterization, only $t$ and $-t$ produce the same value (and these never coincide as $t\ne 0$), so we actually have 
 \[T:=\prod_{a\in A}=(-1)^{\frac{p+1}{4}}\prod_{\substack{t\in\mathbb{F}_p[i]\\t^{\frac{p+1}{2}}-i=0}}(t+1/t).\] This becomes 
 \[T=(-1)^{\frac{p+1}{4}}\prod_{\substack{t\in\mathbb{F}_p[i]\\t^{\frac{p+1}{2}}-i=0}}\frac{(t+i)(t-i)}{t}.\] Letting $f(x)=x^{\frac{p+1}{2}}-i$, we see that this product becomes 
 \[T=(-1)^{\frac{p+1}{4}}\frac{f(-i)f(i)}{f(0)}.\] It is tedious but straightforward to see that this always gives $2$.

 \textbf{Case 2}: $p\equiv 1\pmod{4}$. Let $\ell$ be some QNR mod $p$ (sadly we don't have any free QNRs like $-1$ is for $3\pmod{4}$ primes). Let $\omega$ be a formal square root of $\ell$, and work in $\mathbb{F}_{p^2} = \mathbb{F}_p[\omega]$. Then, if $a\in A$, we have 
 \[\omega^2 a = x^2\quad\text{and}\quad\omega^2(4-a)=y^2,\] so $x^2+y^2=4\omega^2$, so factoring and doing the same as before, we get 
 \[x=\omega(t+1/t)\quad\text{and}\quad y=\frac{\omega}{i}(t-1/t)\] for some $t\in\mathbb{F}_p[\omega]$. These are in $\mathbb{F}_p$ if and only if $t^{p-1}+1=0$ (using frobenius). We see $a=(t+1/t)^2$, so 
 \[A = \{(t+1/t)^2 : t\in\mathbb{F}_p[\omega], t^{p-1}+1=0\}.\] Doing a similar analysis as before, the desired product becomes 
 \[T:=\prod_{a\in A}=(-1)^{\frac{p-1}{4}}\prod_{\substack{t\in\mathbb{F}_p[\omega]\\t^{\frac{p-1}{2}}-i=0}}\frac{(t+i)(t-i)}{t},\] and letting $f(x)=x^{\frac{p-1}{2}}-i$, we get 
 \[T=(-1)^{\frac{p-1}{4}}\frac{f(-i)f(i)}{f(0)},\] which again is always $2$.

 So the answer is $\boxed{2}$.

 


	\item (\textit{USAMO 2020}) Suppose that $(a_1, b_1), (a_2, b_2), \ldots , (a_{100}, b_{100})$ are distinct ordered pairs of nonnegative integers. Let $N$ denote the number of pairs of integers $(i, j)$ satisfying $1 \le i < j \le 100$ and $|a_ib_j - a_j b_i|=1$. Determine the largest possible value of $N$ over all possible choices of the $100$ ordered pairs.


 



 



\textbf{Solution}

Let's start off with just $(a_1, b_1), (a_2, b_2)$ and suppose that it satisfies the given condition. We could use $(1, 1), (1, 2)$ for example. We should maximize the number of conditions that the third pair satisfies. We find out that the third pair should equal $(a_1+a_2, b_1+b_2)$:

 We know this must be true:
 \[|a_1b_2-a_2b_1| = 1\]

 So
 \[a_1b_2-a_2b_1 = 1\]

 We require the maximum conditions for $(a_3, b_3)$
 \[|a_3b_2-a_2b_3| = 1\]
 \[|a_3b_1-a_1b_3| = 1\]

 Then one case can be:
 \[a_3b_2-a_2b_3 = 1\]
 \[a_3b_1-a_1b_3 = -1\]

 We try to do some stuff such as solving for $a_3$ with manipulations:
 \[a_3b_2a_1-a_2b_3a_1 = a_1\]
 \[a_3b_1a_2-a_1b_3a_2 = -a_2\]
 \[a_3(a_1b_2-a_2b_1) = a_1+a_2\]
 \[a_3 = a_1+a_2\]
 \[a_3b_2b_1-a_2b_3b_1 = b_1\]
 \[a_3b_1b_2-a_1b_3b_2 = -b_2\]
 \[b_3(a_1b_2-a_2b_1) = b_1+b_2\]
 \[b_3 = b_1+b_2\]

 We showed that 3 pairs are a complete graph; however, 4 pairs are not a complete graph. We will now show that:
 \[a_4 = a_1+2a_2\]
 \[b_4 = b_1+2b_2\]
 \[|a_1b_1+2a_2b_1-a_1b_1-2a_1b_2| = 1\]
 \[2|a_2b_1-a_1b_2| = 1\]

 This is clearly impossible because $1$ is not even and also $|a_2b_1-a_1b_2| = 1$.The answer is as follows:
 \[0+1+2+\ldots+2\]$a_1$ has $0$ subtractions that follow condition while $a_2$ has $1$ and then the rest has $2$.There are $n$ terms, so our answer be $2n-3$ and in case of $n=100$ that means 
 \[\boxed{N=197}.\]~Lopkiloinm

 



\textbf{Solution 2}

We claim the answer is $197$. 

 Study the points $(0, 0), (a_i, b_i), (a_j, b_j)$. If we let these be the vertices of a triangle, applying shoelace theorem gives us an area of $\frac{1}{2}|0\times{b_i}+{a_i}\times{b_j}+{b_i}\times{0}-0\times{a_i}-{b_i}\times{a_j}-{b_j}\times{0} = \frac{1}{2}|a_ib_j - a_j b_i| = \frac{1}{2}$. Therefore, the triangle formed by the points $(0, 0), (a_i, b_i), (a_j, b_j)$ must have an area of $\frac{1}{2}$. 

 Two cases follow.Case 1: Both $(a_i, b_i), (a_j, b_j)$ have exactly one coordinate equal to $0$. Here, one point must be on the $x$ axis and the other on the $y$ axis in order for the triangle to have a positive area. For the area of the triangle to be $\frac{1}{2}$, it follows that the points must be $(1, 0), (0, 1)$ in some order. 

 Case 2: At least one of $(a_i, b_i), (a_j, b_j)$ does not have exactly one coordinate equal to $0$. Define $S[l]$ to be a list of lines such that each line in the list has some two lattice points that, with $(0, 0)$, form a triangle with area $\frac{1}{2}$. Note that for any such line that passes through such two lattice points, we may trivially generate infinite lattice points on the line that have nonnegative coordinates. 

 Note that lines $y=1$ and $x=1$ are included in $S[l]$, because the points $(1, 1), (2, 1)$ serve as examples for $y=1$ and $(1, 1), (1, 2)$ serve as examples for $x=1$. For the optimal construction, include the points $(1, 0)$ and all the points $(0, 1), (0, 2), (0, 3), ... , (0, 99)$, in that order. In this case, every adjacent pair of points would count ($98$), as well as picking $(0, 1)$ and a nonadjacent point ($99$), so this would be $98+99=197$. 

 To prove that this is the maximum, consider the case where some $n$ number of points were neither on $x=1$ nor on $y=1$. In this case, we would be removing $n$ adjacent pairs and $n$ options to choose from after choosing $(0, 0)$, resulting in a net loss of $2n$. By having $n$ points on some other combination of lines in $S[l]$, we would trivially have a maximum gain of $n-1$ pairs of points on the lines such that there are no lattice points between those pairs. Because these points are not on $x=1$ or $y=1$, the altitude from a given point to the line formed by $(0, 0)$ and $(0, 1)$ and $(1, 0)$ is not $1$, and so the area of the triangle cannot be $\frac{1}{2}$. Thus, by not having all points on lines $x=1$ and $y=1$, we cannot exceed the maximum of $197$. Thus, $\boxed{197}$ is our answer. 

 


 


	\item (\textit{USAMO 2020}) A finite set $S$ of points in the coordinate plane is called \textit{overdetermined} if $|S| \ge 2$ and there exists a nonzero polynomial $P(t)$, with real coefficients and of degree at most $|S| - 2$, satisfying $P(x) = y$ for every point $(x, y) \in S$.


 For each integer $n \ge 2$, find the largest integer $k$ (in terms of $n$) such that there exists a set of $n$ distinct points that is not overdetermined, but has $k$ overdetermined subsets.


 



\textbf{Solution}

The answer is $2^{n-1}-n$. To construct this, have $n-1$ points on a horizontal line and $1$ point not on it. Then, any subset that does not include the last point is overdetermined.\\

 For the bound, the main idea of the problem is the following claim.\\

 Claim: If $S$ and $T$ are overdetermined subsets with $|S|=|T|=k$ but $S$ and $T$ only differ by one element (one is swapped out for another), then $S\cup T$ is overdetermined.\\

 Consider the minimal polynomial of $S$. It has degree at most $k-2$ since $S$ is overdetermined, and similarly with the minimal polynomial of $T$. However, there is a unique polynomial of degree at most $k-2$ passing through $S\cap T$ since it has $k-1$ points. Thus, this unique polynomial also passes through $S$ and $T$. Thus, this polynomial passes through all points of $S\cup T$, as desired.\\

 We now use this claim to inductively bound the number of overdetermined subsets with $k$ elements. Let $m_k$ denote the number of overdetermined subsets of size $k$. Clearly, $m_n=0$.\\

 Claim 2: We have 
 \[m_k\leq \frac{m_{k+1}(k+1)+{n\choose k+1}-m_{k+1}}{n-k}.\] Proof: Each overdetermined subset of size $k+1$ can have up to $k+1$ overdetermined subsets of size $k$, but by the previous claim each non-overdetermined subset of size $k+1$, which there are ${n\choose k+1}-m_{k+1}$ of, can have at most one. Each subset of size $k$ feeds into $n-k$ subsets of size $k+1$. Note that this is increasing in $m_{k+1}$.\\

 Claim 3: We have 
 \[m_k\leq {n-1\choose k}.\] Proof: Use induction. Note that the expression in Claim 2 is nondecreasing in $m_{k+1}$, and verify that 
 \[\frac{{n-1\choose k+1}(k+1)+{n\choose k+1}-{n-1\choose k+1}}{n-k}={n-1\choose k}.\]

 Thus, we have that the number of overdetermined subsets is at most 
 \[\sum_{k=2}^n {n-1\choose k}=2^{n-1}-n,\] as desired.

 


 


	\item (\textit{USAMO 2020}) Let $n \ge 2$ be an integer. Let $x_1 \ge x_2 \ge \cdots \ge x_n$ and $y_1 \ge y_2 \ge \cdots \ge y_n$ be $2n$ real numbers such that

 \begin{align*} 0 &= x_1 + x_2 + \cdots + x_n = y_1 + y_2 + \cdots + y_n\\ \text{and  }1 &= x_1^2+x_2^2+\cdots+x_n^2=y_1^2+y_2^2+\cdots+y_n^2. \end{align*}

Prove that

 \[\sum_{i=1}^n(x_iy_i-x_iy_{n+1-i})\ge\frac{2}{\sqrt{n-1}}.\]


 



\textbf{Solution}

Let $u_i = x_i - x_{n+1-i}$ and $v_i = y_i - y_{n+1-i}$ for $1 \le i \le n$. 

 Since $x_1 \ge x_2 \ge \cdots \ge x_n$ and $y_1 \ge y_2 \ge \cdots \ge y_n$, both $(u_i)$ and $(v_i)$ are non-increasing sequences satisfying $u_{n+1-i} = -u_i$ and $v_{n+1-i} = -v_i$.

 We can express the target sum in terms of $u_i$ and $v_i$:
 \[\sum_{i=1}^n (x_i y_i - x_i y_{n+1-i}) = \sum_{i=1}^n x_i v_i = \frac{1}{2} \sum_{i=1}^n u_i v_i\]

 By the Rearrangement Inequality, because $(y_i)$ is non-increasing and $(y_{n+1-i})$ is non-decreasing:
 \[\sum_{i=1}^n y_i y_{n+1-i} \le \frac{1}{n}\left(\sum_{i=1}^n y_i\right)^2 = 0\]

 This gives a lower bound on the sum of squares of $v_i$:
 \[\sum_{i=1}^n v_i^2 = 2\sum_{i=1}^n y_i^2 - 2\sum_{i=1}^n y_i y_{n+1-i} \ge 2\]

 Similarly, $\sum_{i=1}^n u_i^2 \ge 2$.

 Applying summation by parts along with the Cauchy-Schwarz inequality for ordered sequences subject to $\sum x_i = \sum y_i = 0$ and $\sum x_i^2 = \sum y_i^2 = 1$ yields:
 \[\frac{1}{2} \sum_{i=1}^n u_i v_i \ge \frac{2}{\sqrt{n-1}}\]

 Equality is achieved when $x_1 = y_1 = \frac{1}{\sqrt{2}}$, $x_n = y_n = -\frac{1}{\sqrt{2}}$, and all remaining middle variables are equal to $0$.

 


\end{enumerate}